
       \documentclass[fleqn]{article}      
        \usepackage{amsmath}
        \usepackage{fontspec}
        \usepackage{kotex} % 한국어 지원  

        \begin{document}      
        

쉬움: \\\\  
1. 적분의 정의 - 함수의 면적을 구하는 수학적 방법입니다. \\\\  
2. 정적분의 정의 - 구간 [a, b]에서 함수의 면적을 구하는 것입니다. \\\\  
3. \( \int x \, dx = \frac{x^2}{2} + C \) - 부정적분의 결과입니다. \\\\  
4. \( \int 2x \, dx = x^2 + C \) - 부정적분의 결과입니다. \\\\  
5. \( \int_0^1 x^2 \, dx = \frac{1}{3} \) - 면적을 구한 결과입니다. \\\\

중간: \\\\  
1. \( \int (3x^2 + 2x + 1) \, dx = x^3 + x^2 + x + C \) - 각 항을 적분한 결과입니다. \\\\  
2. 적분의 기본 정리 - 미분과 적분은 서로 역 과정입니다. \\\\  
3. \( \int \sin(x) \, dx = -\cos(x) + C \) - 부정적분의 결과입니다. \\\\  
4. \( \int_1^2 (4x - 1) \, dx = 6 \) - 면적을 구한 결과입니다. \\\\  
5. \( \int e^x \, dx = e^x + C \) - 부정적분의 결과입니다. \\\\

어려움: \\\\  
1. \( \int (x^2 \sin(x)) \, dx \) - 부분적분을 사용하면 \( -x^2 \cos(x) + 2 \int x \cos(x) \, dx \)가 됩니다. \\\\  
2. \( \int_0^{\pi} \sin^2(x) \, dx = \frac{\pi}{2} \) - 삼각함수의 적분을 이용한 결과입니다. \\\\  
3. \( \int \frac{1}{x \ln(x)} \, dx = \ln(\ln(x)) + C \) - 적분의 결과입니다. \\\\  
4. \( \int_0^2 (x^3 - 4x + 6) \, dx = \frac{16}{3} \) - 면적을 구한 결과입니다. \\\\  
5. \( \int (2x^3 - 3x^2 + 4) \, dx = \frac{1}{2} x^4 - x^3 + 4x + C \) - 부정적분의 결과입니다. \\\\      
        \end{document} 
    